\section{Conclusion}

Ola por temrinar


\subsection{Future Work}

Modeling poblational dynamic for depreators, instead of assuming a constant density, would allow for a more realistic relationship and observe oscillations like those of the Lotka-Volterra model . It is also interesting to analyze under Kolmogorov's theorem, which allows saturating and modulating the growth of prey and predators to obtain oscillations without losing realism, as discussed in chapter 4 of Nowak and in chapter 6 of Brauer\cite{nowakEvolutionaryDynamicsExploring2006, brauerMathematicalModelsPopulation2012}. It would also be valuable to incorporate evolutionary game theory on graphs or spatially structured games (chapters 8 and 9 of Nowak), to better understand why animals like guppies form shoals under predation and how they behave within the group; although cellular automata-type simulation could be costly due to its execution time, kaleidoscopic graphs like those proposed by Nowak remain a potentially useful tool for understanding how the fitness landscape is influenced \cite{nowakEvolutionaryDynamicsExploring2006}. For this aport could be useful the thesis of Seghers to understand the behaviour of guppies according the space \cite{hendrikANALYSISGEOGRAPHICVARIATION, reznickExperimentalStudiesEvolution2019}. \\